\chapter{Disegno sperimentale canonico}\label{chap:experiments}
\section{Sette celle e dodici voci}
Le sette celle sono due baseline eval-only (base zero/few), SFT zero-shot e il
$2\times2$ GRPO: inizializzazione base/SFT per prompt train zero/few. La coda ha
12 voci: 2 eval baseline più train+eval per le 5 celle apprese. Ogni eval appresa
esegue zero-shot e retrieval. Non contare le due gambe come metodi diversi.

\begin{center}\begin{tabular}{lll}\toprule
Metodo & Inizializzazione & Prompt train\\\midrule
SFT & base & zero-shot\\
GRPO & base & zero-shot\\
GRPO & base & few-shot\\
SFT-GRPO & SFT & zero-shot\\
SFT-GRPO & SFT & few-shot\\\bottomrule
\end{tabular}\end{center}

\section{Config e artifact}
La radice è \pathcode{experiments/configs/qwen25-05b/}; \pathcode{base.yaml} è
ricetta condivisa, non cella. I runnable dichiarano \pathcode{model_tag},
\pathcode{method}, \pathcode{train_prompt_mode}, \pathcode{variant} e
\pathcode{kind}. Checkpoint/log seguono
\begin{lstlisting}
experiments/{checkpoints,logs}/qwen25-05b/<method>/<train-prompt>/[ablations/<variant>/]run_<timestamp>/
\end{lstlisting}
e risultati
\begin{lstlisting}
experiments/results/qwen25-05b/<method>/<train-prompt>/[ablations/<variant>/]eval-<prompt>/run_<timestamp>/
\end{lstlisting}
Le baseline usano \pathcode{baseline/<prompt>/run_*}. Identità config e artifact
devono concordare; non copiare YAML in gerarchie parallele.

\section{Ablazioni fuori matrice}
PDA, hot e cinque reward sono manuali. Rollout/reward/Markov sono probe su
generazioni congelate. Structured-loss è un benchmark head-only. Nessuno entra
nella campagna primaria senza nuova preregistrazione.
